\documentclass[11pt]{article}
\usepackage{amsmath,amssymb,amsthm,geometry,hyperref}
\geometry{margin=1in}
\newtheorem{theorem}{Theorem}
\newtheorem{lemma}{Lemma}
\newcommand{\ket}[1]{\lvert #1\rangle}
\DeclareMathOperator{\atanTwo}{atan2}
\title{Preparing a Hermite-smoothed quantum initial state}
\author{ASPBE / QuantumComputinglib}
\date{}
\begin{document}
\maketitle

\section{The state, before the circuit}
Fix a nonnegative integer $k$, a register width $n\geq1$, and a truncation
parameter $L>0$. Put $N=2^n$ and $p_j=-\pi L+2\pi Lj/N$ for $0\leq j<N$.
The integer label uses little-endian weights:
$j=\sum_{r=0}^{n-1}2^rq_r$. Thus $q_0$ is the least significant bit,
while the printed ket is $\ket{q_{n-1}\cdots q_0}$.
The required function and quantum state are
\begin{align*}
g_k(p)&=\begin{cases}
 e^p,&p<-1,\\
 P_k(p),&-1\leq p\leq0,\\
 e^{-p},&p>0,
\end{cases}\\
Z_k&=\sum_{j=0}^{N-1}g_k(p_j)^2,
&\ket{g_k}_p&=Z_k^{-1/2}\sum_{j=0}^{N-1}g_k(p_j)\ket j.
\end{align*}
Here $P_k$ has degree at most $2k+1$ and satisfies
\[
P_k^{(a)}(-1)=e^{-1},\qquad P_k^{(a)}(0)=(-1)^a
\quad(0\leq a\leq k).
\]
These are the Hermite conditions in Jin--Liu--Ma,
\emph{Schr\"odingerisation based computationally stable algorithms for
ill-posed problems in partial differential equations},
\href{https://arxiv.org/html/2403.19123v3#S4.SS3}{arXiv:2403.19123v3,
Section 4.3, equations (4.31)--(4.32)}.
The amplitudes are $g_k(p_j)$, not their square roots. The probability of
observing label $j$ is $g_k(p_j)^2/Z_k$.

\section{An explicit polynomial, with no interpolation oracle}
For $0\leq r\leq k$, define
\[
a_{k,r}=\sum_{m=0}^{r}\frac{\binom{k+r-m}{k}}{m!},
\qquad A_k(t)=\sum_{r=0}^{k}a_{k,r}t^r.
\]
With $t=p+1$, set
\[
P_k(p)=e^{-1}(1-t)^{k+1}A_k(t)+t^{k+1}A_k(1-t).
\]
This explicit coefficient formula is the construction used by ASPBE; the
source paper specifies the endpoint conditions.

\begin{lemma}[Endpoint jets and uniqueness]
The displayed polynomial has degree at most $2k+1$ and satisfies all the
Hermite conditions above. No other polynomial of that degree does so.
\end{lemma}
\begin{proof}
The binomial series gives
$(1-t)^{-(k+1)}=\sum_{s\geq0}\binom{k+s}{k}t^s$.
The coefficient of $t^r$ in $e^t(1-t)^{-(k+1)}$ is therefore $a_{k,r}$.
Consequently $(1-t)^{k+1}A_k(t)$ and $e^t$ have identical coefficients
through degree $k$. The other summand contains the factor $t^{k+1}$,
so its first $k$ derivatives vanish at $t=0$. Multiplication by $e^{-1}$
gives the left endpoint values. At $t=1$, exchange $t$ and $1-t$;
each derivative contributes the appropriate factor $(-1)^a$.
Each product has degree at most $2k+1$.
If two interpolants exist, their difference has a zero of multiplicity
at least $k+1$ at both endpoints. It is divisible by
$p^{k+1}(p+1)^{k+1}$, of degree $2k+2$, so the difference is zero.
\end{proof}

\begin{lemma}[Positivity]
The function $g_k$ is positive on the real line; in particular $Z_k>0$.
\end{lemma}
\begin{proof}
All $a_{k,r}$ are positive. For $0<t<1$, both summands in $P_k$ are
positive. At the endpoints their values are $e^{-1}$ and $1$.
The two exterior exponentials are positive as well.
\end{proof}

For $k=1$ the formula reduces to
\[
P_1(p)=(-3+3e^{-1})p^3+(-5+4e^{-1})p^2-p+1.
\]
Its second derivative at zero is $-10+8e^{-1}\ne1$.
Thus the piecewise cubic is $C^1$, not $C^2$. The $C^2$ statement printed
after source equation (6.1) is not adopted here. This case certifies state
preparation; it does not certify the source paper's PDE error estimates.

\section{Construct the rotations from subtree masses}
Write $f_j=g_k(p_j)$. For a depth $d$ and low-bit label $s$, define
\[
m_{d,s}=\sum_{j\bmod 2^d=s}f_j^2,
\quad 0\leq d\leq n,\quad 0\leq s<2^d.
\]
The root mass is $m_{0,0}=Z_k$ and the leaves have $m_{n,j}=f_j^2$.
Partitioning each subtree by its next bit proves
\[
m_{d,s}=m_{d+1,s}+m_{d+1,s+2^d}.
\]
At layer $d$, controlled on $(q_0,\ldots,q_{d-1})$ representing $s$,
apply to $q_d$ the gate
\[
R_y(\theta_{d,s}),\qquad
\theta_{d,s}=2\atanTwo\bigl(\sqrt{m_{d+1,s+2^d}},
                                  \sqrt{m_{d+1,s}}\bigr).
\]
The convention is $R_y(\theta)=
\begin{pmatrix}\cos(\theta/2)&-\sin(\theta/2)\\
\sin(\theta/2)&\cos(\theta/2)\end{pmatrix}$.
For a zero-mass subtree the angle is zero. Its incoming amplitude is already
zero, so this choice does not alter the target.

\begin{theorem}[Deterministic preparation]
The product of these uniformly controlled rotations is unitary and maps
$\ket{0^n}$ to $\ket{g_k}_p$.
\end{theorem}
\begin{proof}
Each fixed control word selects a two-dimensional real rotation. Such a
rotation is orthogonal, hence unitary over the complex numbers. The selected
blocks act on orthogonal control subspaces, and their product is unitary.
After depth $d$ the state is
\[
\sum_{s=0}^{2^d-1}\sqrt{m_{d,s}/Z_k}\,
   \ket{0^{n-d}}\ket{s}.
\]
This holds initially. A rotation splits the amplitude of a parent according
to the square roots of its two child masses, proving the next depth.
At depth $n$, positivity gives $\sqrt{f_j^2}=f_j$, which is the claimed
normalized amplitude. There is no measurement or postselection.
\end{proof}

\section{Global smoothness and exact formal scope}
The endpoint jets also certify that the piecewise function belongs to
$C^k(\mathbb R)$. A generic gluing argument proceeds by induction on the
derivative order: equality of endpoint values gives continuity; equality
of the one-sided derivatives gives differentiability at the junction by
the slope-limit characterization. Applying the induction to the derivative
pieces gives the claim at both junctions.
This is the Lean theorem
\texttt{HermiteSmoothness.smoothInitial\_contDiff}.
The state-preparation root
\texttt{HermiteStatePreparation.hermiteStatePreparation\_complete}
separately certifies normalization, primitive-circuit action, unitarity and
gate counts. Neither root asserts a PDE discretization error, a Sobolev
estimate, or an error bound for machine-rounded angles.

\section{Elementary gates and resource accounting}
A uniformly controlled rotation is compiled recursively using half-sums
and half-differences of its angle table, separated by a CNOT and followed
by a second CNOT. For a control bit $b$, the identity
$X R_y(\beta)X=R_y(-\beta)$ makes the selected angle either
the sum or the difference of the two half-angles.
Induction on the number of controls proves equality of the full matrices.
The reference compiler has $2^d$ rotations and $2(2^d-1)$ CNOTs for $d$
controls. Summing over $d=0,\ldots,n-1$ gives
\[
\#R_y=N-1,\qquad \#\mathrm{CNOT}=2(N-1-n),
\qquad G=3(N-1)-2n.
\]
No auxiliary qubits or state-preparation oracle calls are used. These
counts concern this reference compiler, not an optimality theorem.
Classical coefficient evaluation, sample evaluation, mass reduction and
angle calculation are part of the construction cost. Floating-point
evaluation and angle synthesis are separate from ideal real-angle gates.

\section{Reproduce and inspect the evidence}
The case page links complete Lean modules, the executable generator,
OpenQASM output, a gate diagram, and the numerical acceptance report.
The numerical report records independent NumPy, Qiskit and OpenQASM replay
for finite parameter choices. Those tolerances are not substitutes for
the symbolic Lean state-action theorem. Likewise, the general exact gate
theorem does not assert that machine-rounded rotation angles have zero error.
Every public compilation claim is checked against the repository's current
Lean build before deployment.
\end{document}
